% ------------------------------------------------------------------
% Proposed analytical section: null transport and kernel lumpability
% Intended to be inserted after the current native/index-resolution paths.
% Requires: \usepackage{amsthm}
% Suggested theorem declarations in preamble:
% \newtheorem{proposition}{Proposition}
% \newtheorem{corollary}{Corollary}
% \theoremstyle{definition}\newtheorem{definition}{Definition}
% \theoremstyle{remark}\newtheorem{remark}{Remark}
% ------------------------------------------------------------------

\subsection{Null transport and kernel lumpability}\label{null-transport-lumpability}

Let $(\mathsf X,\mathcal A)$ denote the native-resolution trajectory space and
$(\mathsf Y,\mathcal B)$ a derived-data space.  Let
\[
\mathcal G:\mathsf X\to\mathsf Y
\]
be a deterministic preprocessing map.  In the main application,
$\mathcal G=\mathcal P=D\circ\mathcal O$; when isolating thresholding and
aggregation, we take $\mathcal G=\mathcal O$.  A native-resolution surrogate
procedure is represented by a Markov kernel $K_X(x,dx^\star)$ on $\mathsf X$.
For an observed native trajectory $x$, define its \emph{transported native
null} on the derived space by
\begin{equation}
L_x(B)
:=K_X\!\left(x,\mathcal G^{-1}(B)\right),
\qquad B\in\mathcal B.
\label{eq:transported-native-null}
\end{equation}
Thus $L_x=\mathcal G_\#K_X(x,\cdot)$ is the law of
$\mathcal G(X^\star)$ when $X^\star\sim K_X(x,\cdot)$.

\begin{definition}[Strong $\mathcal G$-lumpability of a surrogate kernel]
The native kernel $K_X$ is strongly $\mathcal G$-lumpable if
\begin{equation}
L_x=L_{x'}
\qquad\text{whenever}\qquad
\mathcal G(x)=\mathcal G(x').
\label{eq:strong-G-lumpability}
\end{equation}
Equivalently, the transported null depends on the native trajectory only
through the derived observation $y=\mathcal G(x)$.
\end{definition}

This is lumpability of the \emph{randomization kernel}, not a claim that the
physical climate process itself is a lumpable Markov chain.  In a finite
state space, Eq.~\eqref{eq:strong-G-lumpability} reduces to the classical
Kemeny--Snell block-transition criterion.

\begin{proposition}[Quotient-null criterion]\label{prop:quotient-null}
Equip $\mathsf Y=\mathcal G(\mathsf X)$ with the quotient sigma-algebra induced
by $\mathcal G$.  The following are equivalent.
\begin{enumerate}
\item $K_X$ is strongly $\mathcal G$-lumpable.
\item There exists a Markov kernel $\bar K_Y$ on $\mathsf Y$ such that
\begin{equation}
\mathcal G_\#K_X(x,\cdot)
=\bar K_Y(\mathcal Gx,\cdot)
\qquad\text{for every }x\in\mathsf X.
\label{eq:quotient-kernel}
\end{equation}
\end{enumerate}
On $\mathcal G(\mathsf X)$ the kernel $\bar K_Y$ is unique.
\end{proposition}

\begin{proof}
If (2) holds and $\mathcal Gx=\mathcal Gx'$, both transported laws equal
$\bar K_Y(\mathcal Gx,\cdot)$, proving (1).  Conversely, if (1) holds, define
$\bar K_Y(y,B)=L_x(B)$ for any $x\in\mathcal G^{-1}(\{y\})$.  Condition
\eqref{eq:strong-G-lumpability} makes this definition independent of the
representative $x$; measurability follows from the quotient sigma-algebra.
Uniqueness on the image of $\mathcal G$ is immediate.
\end{proof}

Proposition~\ref{prop:quotient-null} separates two questions that are distinct
in surrogate inference.  First, does a derived-space null exist that exactly
represents the native randomization after preprocessing?  Second, if it does,
is the chosen index-level surrogate kernel $K_Y$ equal to that quotient
kernel?  The first is a lumpability question; the second is a surrogate-contract
question.

\subsection{A null-transport defect and an unavoidable-error bound}

For $y\in\mathsf Y$, write
$\mathcal F_y=\{x\in\mathsf X:\mathcal Gx=y\}$ for the fiber over $y$ and define
its total-variation transport diameter by
\begin{equation}
\Delta_L(y)
:=\sup_{x,x'\in\mathcal F_y}
 d_{\mathrm{TV}}(L_x,L_{x'}).
\label{eq:lumpability-defect}
\end{equation}
Then $\Delta_L(y)=0$ for every $y$ if and only if $K_X$ is strongly
$\mathcal G$-lumpable.

\begin{proposition}[Unavoidable error of an aggregate-only null]
\label{prop:half-diameter}
Let $Q_Y(y,\cdot)$ be any null kernel that has access only to the derived
observation $y$.  For every fiber,
\begin{equation}
\sup_{x\in\mathcal F_y}
 d_{\mathrm{TV}}\!\left(L_x,Q_Y(y,\cdot)\right)
\ge \frac{1}{2}\Delta_L(y).
\label{eq:half-diameter-bound}
\end{equation}
Consequently, if $\Delta_L(y)>0$, no aggregate-only null can reproduce the
native transported null uniformly over that fiber.
\end{proposition}

\begin{proof}
For any $x,x'\in\mathcal F_y$, the triangle inequality gives
\[
d_{\mathrm{TV}}(L_x,L_{x'})
\le d_{\mathrm{TV}}(L_x,Q_Y(y,\cdot))
+d_{\mathrm{TV}}(Q_Y(y,\cdot),L_{x'}).
\]
At least one term on the right is therefore at least half the left-hand side.
Taking the supremum over $x,x'$ proves Eq.~\eqref{eq:half-diameter-bound}.
\end{proof}

This result gives an operational meaning to failure of lumpability.  If the
native null uses microstate information that $\mathcal G$ removes, then the
aggregate observation alone does not identify a unique transported null.

\begin{corollary}[Projection to a statistic]\label{cor:statistic-projection}
Let $T:\mathsf Y\to\mathsf Z$ be any measurable statistic.  Then
\begin{equation}
d_{\mathrm{TV}}(T_\#L_x,T_\#L_{x'})
\le d_{\mathrm{TV}}(L_x,L_{x'}).
\label{eq:data-processing-TV}
\end{equation}
For any rejection set $R\subseteq\mathsf Z$,
\begin{equation}
\left|
L_x\{T\in R\}-L_{x'}\{T\in R\}
\right|
\le d_{\mathrm{TV}}(T_\#L_x,T_\#L_{x'}).
\label{eq:decision-bound}
\end{equation}
Thus kernel non-lumpability can be invisible to a particular statistic, while
a difference in a statistic-level null distribution is a direct witness of
kernel-level non-lumpability.
\end{corollary}

\begin{proof}
Equation~\eqref{eq:data-processing-TV} is the data-processing inequality for
total variation under a measurable map.  Equation~\eqref{eq:decision-bound}
follows from the definition of total variation.
\end{proof}

\subsection{Model-dependent transport when strong lumpability fails}

Strong lumpability is sufficient for a unique quotient null, but it need not
hold.  Suppose a benchmark null model $P_0$ on $\mathsf X$ admits a regular
conditional distribution $\mu_y(dx)=P_0(dx\mid \mathcal G(X)=y)$.  Define
\begin{equation}
\bar K_{P_0}(y,B)
:=\int_{\mathcal F_y} L_x(B)\,\mu_y(dx).
\label{eq:model-dependent-transport}
\end{equation}

\begin{proposition}[Model-dependent transported null]
\label{prop:model-dependent-transport}
If $X\sim P_0$ and $X^\star\mid X\sim K_X(X,\cdot)$, then
\begin{equation}
\Pr\{\mathcal G(X^\star)\in B\mid \mathcal G(X)=y\}
=\bar K_{P_0}(y,B).
\label{eq:conditional-transport}
\end{equation}
If $K_X$ is strongly $\mathcal G$-lumpable, then
$\bar K_{P_0}=\bar K_Y$ for every $P_0$.  If strong lumpability fails, the
transported null can depend on how $P_0$ distributes unresolved microstates
within the fiber $\mathcal F_y$.
\end{proposition}

\begin{proof}
Equation~\eqref{eq:conditional-transport} follows from the tower property:
conditional on $X=x$, the probability of
$\mathcal G(X^\star)\in B$ is $L_x(B)$; averaging this quantity over
$P_0(dx\mid \mathcal G(X)=y)$ yields Eq.~\eqref{eq:model-dependent-transport}.
Under strong lumpability, $L_x$ is constant on each fiber, so the integral is
independent of $\mu_y$ and equals the quotient kernel from
Proposition~\ref{prop:quotient-null}.
\end{proof}

\subsection{Minimal spectral-surrogate counterexample}\label{spectral-counterexample}

The following finite example isolates the structural obstruction relevant to
marginal-and-spectrum surrogates.  It is not an analytical representation of
IAAFT itself; rather, it is a minimal exactly solvable surrogate kernel that
preserves the empirical marginal and softly matches a spectral feature.

Let
\[
\Omega=\{z\in\{0,1\}^4:\;z_1+z_2+z_3+z_4=2\},
\]
so every state has exactly two event and two non-event values.  Interpret the
four coordinates as two consecutive two-period blocks and define the block-count
coarse-graining
\begin{equation}
\mathcal O(z)=(z_1+z_2,\;z_3+z_4).
\label{eq:toy-O}
\end{equation}
For $\bar z=1/2$, define the centered circular lag-one score
\begin{equation}
R(z)=\sum_{t=1}^{4}(z_t-\bar z)(z_{t+1}-\bar z),
\qquad z_5=z_1.
\label{eq:toy-R}
\end{equation}
With the discrete Fourier transform convention
$\widehat z_k=\sum_{t=0}^{3}(z_{t+1}-\bar z)e^{-2\pi i kt/4}$,
\[
R(z)=\frac{1}{4}\sum_{k=0}^{3}|\widehat z_k|^2
\cos(2\pi k/4),
\]
so $R$ is a linear functional of Fourier power.

For $\tau\ge0$, define the spectral-tilted permutation kernel
\begin{equation}
K_\tau(x,z)
=
\frac{\exp\{-\tau[R(z)-R(x)]^2\}}
{\sum_{w\in\Omega}\exp\{-\tau[R(w)-R(x)]^2\}},
\qquad z\in\Omega.
\label{eq:toy-kernel}
\end{equation}
All outputs have the same empirical marginal as the input.  At $\tau=0$ the
kernel is uniform over all permutations; increasing $\tau$ increasingly favors
permutations whose lag-one spectral score matches that of the observed
microstate.

\begin{proposition}[Spectral conditioning can destroy lumpability]
\label{prop:toy-counterexample}
For the block-count map \eqref{eq:toy-O}, $K_0$ is strongly
$\mathcal O$-lumpable, whereas $K_\tau$ is not strongly
$\mathcal O$-lumpable for every $\tau>0$.

More explicitly, take
\[
x=(0,1,1,0),
\qquad
x'=(1,0,1,0).
\]
Then
\[
\mathcal O(x)=\mathcal O(x')=(1,1),
\qquad
R(x)=0,
\qquad
R(x')=-1.
\]
Writing $q=e^{-\tau}$, the total-variation distance between the two transported
block-count nulls is
\begin{equation}
 d_{\mathrm{TV}}
 \left(
 \mathcal O_\#K_\tau(x,\cdot),
 \mathcal O_\#K_\tau(x',\cdot)
 \right)
 =
 \frac{1-q^2}{2q^2+5q+2}
 >0
 \qquad (\tau>0).
\label{eq:toy-TV}
\end{equation}
It increases from $0$ at $\tau=0$ to $1/2$ as $\tau\to\infty$.
\end{proposition}

\begin{proof}
Among the six elements of $\Omega$, four have $R=0$,
\[
0011,\;0110,\;1001,\;1100,
\]
and two have $R=-1$,
\[
0101,\;1010.
\]
The possible block-count values are $(2,0)$, $(1,1)$, and $(0,2)$.  For the
microstate $x$, whose target score is $0$, the transported distribution is
\begin{equation}
L_x
=
\frac{1}{4+2q}
\left(1,\;2+2q,\;1\right),
\label{eq:toy-Lx}
\end{equation}
whereas for $x'$, whose target score is $-1$,
\begin{equation}
L_{x'}
=
\frac{1}{2+4q}
\left(q,\;2+2q,\;q\right).
\label{eq:toy-Lxp}
\end{equation}
These distributions coincide only when $q=1$, i.e. $\tau=0$.  Direct
calculation gives Eq.~\eqref{eq:toy-TV}.
\end{proof}

The counterexample also gives an exact statistic-level witness.  Let
\begin{equation}
T_{\mathrm{diff}}(y_1,y_2)=(y_2-y_1)^2.
\end{equation}
Then
\[
\Pr_x\{T_{\mathrm{diff}}=4\}
=\frac{1}{2+q},
\qquad
\Pr_{x'}\{T_{\mathrm{diff}}=4\}
=\frac{q}{1+2q},
\]
and therefore
\begin{equation}
 d_{\mathrm{TV}}((T_{\mathrm{diff}})_\#L_x,(T_{\mathrm{diff}})_\#L_{x'})
 =
 \frac{1-q^2}{2q^2+5q+2}.
\label{eq:toy-Tdiff-TV}
\end{equation}
Thus this simple order-sensitive statistic retains the entire fiber-level
transport defect in the example.  This mirrors the role of $T_{\mathrm{diff}}$
in the within-fiber computational diagnostic, where it was more sensitive to
microstate-dependent null transport than the extremal-index statistic.

Finally, Proposition~\ref{prop:half-diameter} implies that, on the fiber
$\mathcal O^{-1}(1,1)$, every aggregate-only null $Q_Y(1,1,\cdot)$ must incur
worst-case total-variation error at least
\[
\frac{1}{2}
\frac{1-e^{-2\tau}}
{2e^{-2\tau}+5e^{-\tau}+2},
\]
which tends to $1/4$ as $\tau\to\infty$.  The obstruction is therefore not
merely that one particular index-level surrogate is poorly chosen: once the
native kernel uses a microstate spectral feature erased by $\mathcal O$, no
single null based only on the aggregate can be uniformly exact over the fiber.

